\chapter{Геометрия носителя}
\label{ch:carrier}

\section{Что уже введено}
\label{sec:carrier-motivation}

Носитель как счётное множество узлов~$\Lambda$, элементарные шаг~$\hl$ и такт~$\htick$,
тактовая скорость $c_0=\hl/\htick$, окрестность одного такта и поле на узле
появились в конце спуска (\cref{sec:planck-floor}).
Абсолютные условия $A_1$--$A_{16}$ (гл.~\ref{ch:axiom-rationale}, \ref{ch:axioms})
ограничивают класс локальных законов на этом носителе
(\cref{fig:carrier-lattice-field,fig:carrier-field-z-spinor,fig:carrier-field-z-law}).

\begin{figure}[htbp]
  \centering
  \begin{tikzpicture}[carrier panel, x=1.05cm, y=0.92cm]
  \foreach \r in {0,...,3} {
    \foreach \c in {0,...,4} {
      \ifnum\r=1
        \ifnum\c=2
          \node[carrier site, minimum size=5pt] at (\c,\r) {};
        \else
          \node[carrier site, minimum size=3pt] at (\c,\r) {};
        \fi
      \else
        \node[carrier site, minimum size=3pt] at (\c,\r) {};
      \fi
    }
  }
  \foreach \r in {0,...,3} \foreach \c in {0,...,3} {
    \draw[carrier muted, line width=0.35pt] (\c,\r) -- (\c+1,\r);
    \draw[carrier muted, line width=0.35pt] (\c,\r) -- (\c,\r+1);
  }
  \CarrierDimH{2}{3}{1}{$\hl$}
  \node[anchor=west, inner sep=2pt] at (2.06,1.04) {$x\in\Lambda$};
  \draw[carrier object] (3.05,-0.42) -- (4.05,-0.42);
  \draw[carrier line] (3.05,-0.47) -- (3.05,-0.37) (4.05,-0.47) -- (4.05,-0.37);
  \node at (3.05,-0.58) {\scriptsize $t$};
  \node at (4.05,-0.58) {\scriptsize $t{+}1$};
  \CarrierDimH{3.05}{4.05}{-0.42}{$\htick$}
  \node[anchor=west, font=\footnotesize] at (0.05,3.35) {$\Lambda$ --- счётное множество узлов};
  \node at (1.55,-0.05) {\footnotesize $c_0=\hl/\htick$};
\end{tikzpicture}

  \caption{Носитель $\Lambda$, шаг $\hl$, такт $\htick$.}
  \label{fig:carrier-lattice-field}
\end{figure}

\begin{figure}[htbp]
  \centering
  \begin{tikzpicture}[carrier panel]
  \filldraw[carrier object] (0,0) circle (2.2pt);
  \node[anchor=east, inner sep=2pt] at (-0.06,0.02) {$x$};
  \node[draw, carrier object, carrier fill, align=center, inner sep=4pt, anchor=west] (zbox) at (0.55,0.15)
    {$z(x)=\binom{z_1}{z_2}\in\mathbb{C}^2$};
  \draw[carrier arrow] (0.06,0.04) -- (zbox.west);
\end{tikzpicture}

  \caption{Состояние узла: спинор $z(x)\in\mathbb{C}^2$.}
  \label{fig:carrier-field-z-spinor}
\end{figure}

\begin{figure}[htbp]
  \centering
  \begin{tikzpicture}[carrier panel]
  \CarrierHexRing{0.36}{0.62}
  \filldraw[carrier object] (0,0) circle (2.2pt);
  \node[anchor=south west, inner sep=1pt] at (-0.1,0.08) {$x$};
  \node[anchor=south, inner sep=2pt] at (0,0.95) {$N(x)$};
  \node[draw, carrier object, carrier fill, minimum width=0.9cm, minimum height=0.9cm] (g) at (2.35,0) {$g$};
  \node[draw, carrier object, circle, minimum size=0.95cm, align=center, inner sep=1pt] (z1) at (4.05,0) {\scriptsize $z(x,t{+}1)$};
  \foreach \a in {-60,0,60} {
    \draw[carrier arrow, carrier muted] ({0.62*cos(\a)}, {0.62*sin(\a)}) -- (g.west);
  }
  \draw[carrier arrow] (g.east) -- (z1.west);
\end{tikzpicture}

  \caption{Локальный закон $g$ на окрестности $N(x)$.}
  \label{fig:carrier-field-z-law}
\end{figure}

\begin{figure}[htbp]
  \centering
  \begin{tikzpicture}[carrier panel, x=0.72cm, y=0.64cm]
  \newcommand{\snapcell}[3]{%
    \fill[gray!#3] (#1,#2) circle (0.17);
    \draw[carrier object] (#1,#2) circle (0.17);
    \node[carrier site, minimum size=2pt] at (#1,#2) {};
  }
  \foreach \r in {0,...,3} \foreach \c in {0,...,3} {
    \draw[carrier muted, line width=0.35pt] (\c,\r) -- (\c+1,\r);
    \draw[carrier muted, line width=0.35pt] (\c,\r) -- (\c,\r+1);
  }
  \snapcell{0}{3}{35}\snapcell{1}{3}{55}\snapcell{2}{3}{35}\snapcell{3}{3}{15}
  \snapcell{0}{2}{55}\snapcell{1}{2}{75}\snapcell{2}{2}{35}\snapcell{3}{2}{55}
  \snapcell{0}{1}{15}\snapcell{1}{1}{55}\snapcell{2}{1}{75}\snapcell{3}{1}{35}
  \snapcell{0}{0}{35}\snapcell{1}{0}{15}\snapcell{2}{0}{55}\snapcell{3}{0}{35}
  \fill[carrier hifill, opacity=0.45] (1,2) circle (0.34);
  \foreach \n in {(1,3),(1,1),(0,2),(2,2)} {
    \fill[carrier hifill, opacity=0.3] \n circle (0.34);
  }
  \draw[carrier object] (1,2) circle (0.28);
  \node[anchor=west] at (1.08,2.08) {$x$};
\end{tikzpicture}

  \caption{Синхронный такт: поле $\Psi^t$ (окрестность чтения выделена).}
  \label{fig:carrier-sync-psi-t}
\end{figure}

\begin{figure}[htbp]
  \centering
  \begin{tikzpicture}[carrier panel, x=0.72cm, y=0.64cm]
  \newcommand{\snapcell}[3]{%
    \fill[gray!#3] (#1,#2) circle (0.17);
    \draw[carrier object] (#1,#2) circle (0.17);
    \node[carrier site, minimum size=2pt] at (#1,#2) {};
  }
  \foreach \r in {0,...,3} \foreach \c in {0,...,3} {
    \draw[carrier muted, line width=0.35pt] (\c,\r) -- (\c+1,\r);
    \draw[carrier muted, line width=0.35pt] (\c,\r) -- (\c,\r+1);
  }
  \snapcell{0}{3}{35}\snapcell{1}{3}{55}\snapcell{2}{3}{55}\snapcell{3}{3}{15}
  \snapcell{0}{2}{75}\snapcell{1}{2}{35}\snapcell{2}{2}{55}\snapcell{3}{2}{55}
  \snapcell{0}{1}{15}\snapcell{1}{1}{75}\snapcell{2}{1}{55}\snapcell{3}{1}{35}
  \snapcell{0}{0}{35}\snapcell{1}{0}{35}\snapcell{2}{0}{55}\snapcell{3}{0}{35}
  \draw[carrier object, dashed] (1,2) circle (0.24);
\end{tikzpicture}

  \caption{Синхронный такт: поле $\Psi^{t+1}$ после параллельного применения $g$.}
  \label{fig:carrier-sync-psi-t1}
\end{figure}

\section{Световая окрестность}
\label{sec:epsilon}

\begin{definition}[Световая $\varepsilon$-окрестность]
\label{def:epsilon}
Для узла $x\in\Lambda$ окрестность одного такта --- узлы, до которых сигнал доходит за~$\htick$:
\[
  N(x)\subset\Lambda.
\]
\end{definition}

\begin{figure}[htbp]
  \centering
  \begin{tikzpicture}[carrier panel, scale=1.35]
  \CarrierHexRing{0.42}{0.72}
  \filldraw[carrier object] (0,0) circle (2.5pt);
  \node[anchor=west, inner sep=2pt] at (0.08,-0.02) {$x$};
  \node[anchor=south, inner sep=3pt] at (0,1.05) {$N(x)$};
  \draw[carrier arrow, carrier muted] (0.42,0) -- (0.72,0);
  \node[anchor=south, carrier muted, inner sep=1pt] at (0.57,-0.12) {$\hl$};
\end{tikzpicture}

  \caption{Световая $\varepsilon$-окрестность $N(x)$: узлы на расстоянии~$\hl$ за один такт~$\htick$.}
  \label{fig:carrier-epsilon}
\end{figure}

\section{Гексагональный срез $(2{+}1)$}
\label{sec:hex-slice}

\subsection{Регулярные замощения плоскости}

Локальный закон на плоскости требует замощения без разрывов.
На плоскости существуют ровно три регулярных типа: треугольник, квадрат и шестиугольник
(\cref{fig:carrier-tiling-triangle,fig:carrier-tiling-square,fig:carrier-tiling-hex}).
Число соседей $|N(x)|$ на одном такте --- отдельный вопрос: узлы стоят в центрах плиток
или в вершинах, и квадратная решётка допускает диагональ $h_L\sqrt{2}$, недостижимую за один~$\htick$
(\cref{fig:carrier-neighbors-tri,fig:carrier-neighbors-sq,fig:carrier-neighbors-hex}).
Наиплотнейшая изотропная упаковка центров даёт гексагональную ячейку Вороного, $|N|=6$
(\cref{fig:carrier-hex-neighbors}).

\begin{figure}[htbp]
  \centering
  \begin{tikzpicture}[carrier panel]
  \foreach \i in {-1,...,1} {
    \foreach \j in {-1,...,1} {
      \pgfmathsetmacro{\x}{\i*0.55 + mod(\j,2)*0.275}
      \pgfmathsetmacro{\y}{\j*0.48}
      \draw[carrier object, carrier fill]
        (\x,\y) -- ++(1,0) -- ++(-0.5,0.866) -- cycle;
    }
  }
  \draw[carrier object, carrier hifill] (0,0) -- (1,0) -- (0.5,0.866) -- cycle;
\end{tikzpicture}

  \caption{Треугольное замощение (выделена одна плитка).}
  \label{fig:carrier-tiling-triangle}
\end{figure}

\begin{figure}[htbp]
  \centering
  \begin{tikzpicture}[carrier panel]
  \foreach \i in {0,...,2} {
    \foreach \j in {0,...,2} {
      \draw[carrier object, carrier fill] (\i,\j) rectangle ++(1,1);
    }
  }
  \draw[carrier object, carrier hifill] (1,1) rectangle ++(1,1);
  \CarrierDimH{1}{2}{-0.35}{$a$}
\end{tikzpicture}

  \caption{Квадратное замощение (выделена одна плитка).}
  \label{fig:carrier-tiling-square}
\end{figure}

\begin{figure}[htbp]
  \centering
  \begin{tikzpicture}[carrier panel]
  \foreach \r in {-1,...,1} {
    \foreach \c in {-1,...,1} {
      \pgfmathsetmacro{\cx}{\c*0.75 + mod(\r,2)*0.375}
      \pgfmathsetmacro{\cy}{\r*0.65}
      \CarrierHex{0.32}{(\cx,\cy)}
    }
  }
  \draw[carrier object, carrier hifill]
    (30:0.32) \foreach \a in {90,150,...,330} { -- (\a:0.32) } -- cycle;
\end{tikzpicture}

  \caption{Шестиугольное замощение (выделена одна плитка).}
  \label{fig:carrier-tiling-hex}
\end{figure}

\begin{figure}[htbp]
  \centering
  \begin{tikzpicture}[carrier panel]
  \foreach \i in {-1,...,1} {
    \foreach \j in {-1,...,1} {
      \pgfmathsetmacro{\x}{\i*0.55 + mod(\j,2)*0.275}
      \pgfmathsetmacro{\y}{\j*0.48}
      \draw[carrier object, carrier fill]
        (\x,\y) -- ++(1,0) -- ++(-0.5,0.866) -- cycle;
    }
  }
  \fill[carrier hifill] (0.5,0.29) circle (2pt);
  \node[carrier site] at (0.5,0.29) {};
\end{tikzpicture}

  \caption{Треугольная решётка: $|N(x)|=3$.}
  \label{fig:carrier-neighbors-tri}
\end{figure}

\begin{figure}[htbp]
  \centering
  \begin{tikzpicture}[carrier panel]
  \foreach \i in {0,...,2} {
    \foreach \j in {0,...,2} {
      \draw[carrier object, carrier fill] (\i,\j) rectangle ++(1,1);
    }
  }
  \foreach \i/\j in {1/1,2/1,0/1,1/2,1/0} {
    \draw[carrier object, carrier hifill] (\i,\j) rectangle ++(1,1);
  }
  \node[carrier site] at (1.5,1.5) {};
  \draw[carrier muted, dashed] (1.5,1.5) -- (2.71,2.71);
  \node[carrier muted, font=\scriptsize] at (2.35,2.45) {$\hl\sqrt{2}$};
\end{tikzpicture}

  \caption{Квадратная решётка: $|N(x)|=4$; диагональ $\hl\sqrt{2}$ недостижима за один такт.}
  \label{fig:carrier-neighbors-sq}
\end{figure}

\begin{figure}[htbp]
  \centering
  \begin{tikzpicture}[carrier panel]
  \CarrierHexNeighborhood{0.36}{0.62}
  \node[carrier site] at (0,0) {};
\end{tikzpicture}

  \caption{Гексагональная решётка: $|N(x)|=6$.}
  \label{fig:carrier-neighbors-hex}
\end{figure}

\begin{figure}[htbp]
  \centering
  \begin{tikzpicture}[carrier panel, scale=1.4]
  \CarrierHexRing{0.4}{0.68}
  \filldraw[carrier object] (0,0) circle (2.5pt);
  \draw[carrier object, carrier muted] (0,0) circle (0.4);
  \draw[carrier arrow] (0,0) -- (0.4,0);
  \node[anchor=south, inner sep=1pt] at (0.2,-0.06) {$a=\hl$};
  \node[anchor=west, inner sep=2pt] at (0.06,0) {$x$};
  \node[anchor=south, inner sep=2pt] at (0,0.92) {$|N(x)|=6$};
\end{tikzpicture}

  \caption{Гексагональная ячейка Вороного: ребро $a=\hl$ и $|N(x)|=6$.}
  \label{fig:carrier-hex-neighbors}
\end{figure}

\subsection{Константы на гексагональном слое}

Центр правильного шестиугольника и его вершина разделены ребром~$a=\hl$.
Апофема --- расстояние от центра до стороны --- есть высота равностороннего треугольника со стороной~$a$:
\[
  R_{\mathrm{in}}=\frac{\sqrt{3}}{2}\,a.
\]
Это чистая геометрия однотактового тела, без соседей
(\cref{fig:carrier-hex-radii}).
Отношение к радиусу до вершины $R_{\mathrm{out}}=a$ даёт геометрический множитель среза
\[
  \kgeom_{\mathrm{hex}}=\frac{R_{\mathrm{in}}}{R_{\mathrm{out}}}=\frac{\sqrt{3}}{2}.
\]

\begin{figure}[htbp]
  \centering
  \begin{tikzpicture}[carrier panel, scale=1.5]
  \def\a{1}
  \pgfmathsetmacro{\rin}{0.866*\a}
  \draw[carrier object] (0,0) circle (\a);
  \draw[carrier muted, dashed] (0,0) circle (\rin);
  \draw[carrier object, carrier fill] (0:\a) \foreach \b in {60,120,...,300} { -- (\b:\a) } -- cycle;
  \node[carrier site] at (0,0) {};
  \draw[carrier arrow] (0,0) -- (\a,0);
  \node[anchor=south, inner sep=1pt] at (0.5*\a,-0.06) {$R_{\mathrm{out}}=a$};
  \draw[carrier arrow, carrier muted] (0,0) -- (0,\rin);
  \node[anchor=west, carrier muted, inner sep=1pt] at (0.08,0.5*\rin) {$R_{\mathrm{in}}$};
  \node[anchor=north, font=\footnotesize] at (0,-1.35)
    {$\kappa_{\mathrm{hex}}=R_{\mathrm{in}}/R_{\mathrm{out}}=\sqrt{3}/2$};
\end{tikzpicture}

  \caption{Вписанный и описанный радиусы правильного шестиугольника: $\kgeom_{\mathrm{hex}}$.}
  \label{fig:carrier-hex-radii}
\end{figure}

Тактовая скорость $c_0=\hl/\htick$ и макроскопическая~$c$ связаны этим же отношением, $c=\kgeom_{\mathrm{hex}} c_0$,
поэтому
\[
  \htick=\kgeom_{\mathrm{hex}}\,t_P,\qquad
  c_0=\frac{2}{\sqrt{3}}\,c,\qquad
  \kappa_{\mathrm{link}}=\frac{1}{|N|}=\frac{1}{6}.
\]
Число~$c$ одно и на срезе, и на трёхмерном носителе, который будет выбран ниже;
меняется только разбиение $c=\kgeom c_0$.

\section{Выбор FCC с двенадцатью соседями}
\label{sec:fcc-choice}

\begin{theorem}[Носитель в $(3{+}1)$]
\label{thm:fcc-choice}
Пусть выполнены: замощение без дыр; единственная длина~$\hl$;
только светоподобные рёбра за один~$\htick$;
изотропия вакуума; один носитель на узел; физическое пространство трёхмерно.
Тогда носитель --- $\Lambda=\mathrm{FCC}$, $|N(x)|=12$ для всех узлов,
$\kappa_{\mathrm{link}}=1/12$.
\end{theorem}

\begin{proof}
Сначала отбираются упаковки с одной длиной ребра и без дыр, затем отсекается вторая координационная оболочка световым конусом.

\emph{Упаковка.}
Ячейку представляем шаром радиуса $\hl/2$: касание двух шаров есть ребро длины~$\hl$.
В трёх измерениях плотнейшая упаковка равных шаров --- гранецентрированная кубическая или гексагональная плотная (теорема Кеплера).
У каждого шара тогда ровно двенадцать соседей на расстоянии~$\hl$.
Простая кубическая решётка имеет шесть соседей, объёмно-центрированная --- восемь; обе реже и отпадают, раз требуется максимальная плотность при одной длине ребра
(\cref{fig:carrier-packing-sc,fig:carrier-packing-bcc,fig:carrier-packing-fcc}).

\begin{figure}[htbp]
  \centering
  \begin{tikzpicture}[carrier panel, iso view, scale=0.95]
  \def\r{0.38}
  \IsoSphereInk{0}{0}{0}{\r}
  \foreach \x/\y/\z in {1,0,0/-1,0,0/0,1,0/0,-1,0/0,0,1/0,0,-1} {
    \IsoSphereMuted{\x}{\y}{\z}{\r}
  }
  \node[anchor=north, font=\footnotesize] at (0,-1.15) {SC, $|N|=6$};
\end{tikzpicture}

  \caption{Простая кубическая упаковка (SC): $|N|=6$.}
  \label{fig:carrier-packing-sc}
\end{figure}

\begin{figure}[htbp]
  \centering
  \begin{tikzpicture}[carrier panel, iso view, scale=0.95]
  \def\r{0.36}
  \IsoSphereInk{0}{0}{0}{\r}
  \foreach \s in {-1,1} {
    \IsoSphereMuted{\s}{\s}{\s}{\r}
    \IsoSphereMuted{\s}{\s}{-\s}{\r}
    \IsoSphereMuted{\s}{-\s}{\s}{\r}
    \IsoSphereMuted{-\s}{\s}{\s}{\r}
  }
  \node[anchor=north, font=\footnotesize] at (0,-1.2) {BCC, $|N|=8$};
\end{tikzpicture}

  \caption{Объёмно-центрированная упаковка (BCC): $|N|=8$.}
  \label{fig:carrier-packing-bcc}
\end{figure}

\begin{figure}[htbp]
  \centering
  \begin{tikzpicture}[carrier panel, iso view, scale=0.95]
  \def\r{0.34}
  \IsoSphereInk{0}{0}{0}{\r}
  \foreach \x/\y/\z in {
    1,1,0/1,-1,0/-1,1,0/-1,-1,0,
    1,0,1/1,0,-1/-1,0,1/-1,0,-1,
    0,1,1/0,1,-1/0,-1,1/0,-1,-1
  } {
    \IsoSphereMuted{\x}{\y}{\z}{\r}
  }
  \node[anchor=north, font=\footnotesize] at (0,-1.25) {FCC, $|N|=12$};
\end{tikzpicture}

  \caption{Гранецентрированная упаковка (FCC): $|N|=12$.}
  \label{fig:carrier-packing-fcc}
\end{figure}

Плотность и число соседей у FCC и HCP одинаковы.
Их различает симметрия.
У FCC есть три равные кубические оси, и осреднение по ним не выделяет направление.
У HCP есть выделенная ось укладки слоёв; она осталась бы в тензоре вторых моментов.
Поэтому носитель --- FCC.

\emph{Световой конус.}
Двенадцать соседей узла в начале координат, при ребре кубической ячейки $\hl\sqrt{2}$, суть все перестановки
\[
  \bigl(\pm\hl/\sqrt{2},\ \pm\hl/\sqrt{2},\ 0\bigr).
\]
Расстояние до начала
\[
  \sqrt{2\cdot(\hl/\sqrt{2})^2}=\hl.
\]
Вторая оболочка --- перестановки $(\pm\hl\sqrt{2},0,0)$, расстояние $\hl\sqrt{2}$.
За один такт $c_0\htick=\hl$, и интервал до узла второй оболочки
\[
  c_0^2\htick^2-(\hl\sqrt{2})^2=\hl^2-2\hl^2=-\hl^2<0.
\]
Интервал пространственноподобен: узел второй оболочки за этот такт недостижим
(\cref{fig:carrier-light-cone-2d,fig:carrier-light-cone-3d}).
Остаётся первая оболочка: $|N(x)|=12$ (\cref{fig:carrier-fcc-shell}).

\begin{figure}[htbp]
  \centering
  \begin{tikzpicture}[carrier panel]
  \CarrierLightCone{1.15}{1.35}
  \fill[carrier muted, opacity=0.25] (-1.35,1.15) -- (0,0) -- (1.35,1.15) -- cycle;
  \node[carrier site] at (0,0) {};
  \node[carrier muted, font=\scriptsize] at (-0.35,0.55) {будущее};
  \draw[carrier object] (0,0) -- (1.0,1.15);
  \node[carrier site, minimum size=2.5pt] at (1.0,1.15) {};
  \draw[carrier muted, dashed] (0,0) -- (1.41,1.15);
  \node[carrier muted, font=\scriptsize] at (1.05,0.75) {$\hl\sqrt{2}$};
  \CarrierDimH{0}{1.0}{1.15}{$\hl$}
  \CarrierDimV{0}{1.15}{0}{$\htick$}
\end{tikzpicture}

  \caption{Световой конус $c_0$ в срезе $(x,t)$: первая оболочка на конусе, вторая ($\hl\sqrt{2}$) --- вне.}
  \label{fig:carrier-light-cone-2d}
\end{figure}

\begin{figure}[htbp]
  \centering
  \begin{tikzpicture}[carrier panel, iso view, scale=0.95]
  \foreach \a in {0,45,...,315} {
    \draw[carrier muted] (0,0,0) -- ({1.1*cos(\a)},{1.1*sin(\a)},1.1);
  }
  \draw[carrier object] (0,0,0) -- (0,0,1.1);
  \fill[carrier fill, opacity=0.5] (0,0,0) -- (1.1,0,1.1) -- (0,1.1,1.1) -- cycle;
  \node[carrier site] at (0,0,0) {};
\end{tikzpicture}

  \caption{Световой конус в объёме $(x,y,t)$.}
  \label{fig:carrier-light-cone-3d}
\end{figure}

\begin{figure}[htbp]
  \centering
  \begin{tikzpicture}[carrier panel, iso view, scale=1.05]
  \def\r{0.36}
  \IsoSphereInk{0}{0}{0}{\r}
  \foreach \x/\y/\z in {
    1,1,0/1,-1,0/-1,1,0/-1,-1,0,
    1,0,1/1,0,-1/-1,0,1/-1,0,-1,
    0,1,1/0,1,-1/0,-1,1/0,-1,-1
  } {
    \IsoSphereMuted{\x}{\y}{\z}{\r}
  }
  \draw[carrier arrow, carrier muted] (0,0,0) -- (1,0,0);
  \node[carrier muted, font=\scriptsize, anchor=west] at (0.55,-0.15) {$\hl$};
\end{tikzpicture}

  \caption{Плотнейшая FCC-упаковка: $|N|=12$ касающихся сфер.}
  \label{fig:carrier-fcc-shell}
\end{figure}

\emph{Связь.}
Изотропия микрошага даёт равные веса на этих двенадцати рёбрах, поэтому
\[
  \kappa_{\mathrm{link}}=1/|N|=1/12.
\]
\end{proof}

\subsection{Геометрия однотактового тела на FCC}
\label{sec:fcc-kappa}

Выпуклая оболочка двенадцати соседей --- кубооктаэдр.
Его вершины --- те же точки $(\pm a/\sqrt{2},\pm a/\sqrt{2},0)$ и перестановки, где $a=\hl$.
Расстояние от центра до вершины уже сосчитано:
\[
  R_{\mathrm{out}}=a.
\]
Четыре точки $(a/\sqrt{2},\pm a/\sqrt{2},0)$ и $(a/\sqrt{2},0,\pm a/\sqrt{2})$ лежат в плоскости $x=a/\sqrt{2}$ и образуют квадрат со стороной~$a$.
Таких квадратных граней шесть.
Расстояние от центра до грани
\[
  R_{\mathrm{in}}^{\mathrm{hull}}=a/\sqrt{2}.
\]
Отношение двух длин одного и того же тела
\[
  \kgeom=R_{\mathrm{in}}^{\mathrm{hull}}/R_{\mathrm{out}}=1/\sqrt{2}
\]
ещё не использует скорость света: это геометрия оболочки.
Треугольная грань через $(a/\sqrt{2},a/\sqrt{2},0)$, $(a/\sqrt{2},0,a/\sqrt{2})$ и $(0,a/\sqrt{2},a/\sqrt{2})$ лежит в плоскости $x+y+z=a\sqrt{2}$ и отстоит от центра на $a\sqrt{2/3}$.
Квадратные и треугольные грани лежат на разных расстояниях, поэтому одной вписанной сферы, касающейся всех граней, нет.
В мост $c=\kgeom c_0$ входит расстояние до квадратной грани
(\cref{fig:carrier-cuboctahedron,fig:carrier-cubocta-meridian}).

\begin{figure}[htbp]
  \centering
  \begin{tikzpicture}[carrier panel, iso view, scale=1.15]
  \foreach \x/\y/\z in {1,1,0/1,-1,0/-1,1,0/-1,-1,0/1,0,1/1,0,-1/-1,0,1/-1,0,-1/0,1,1/0,1,-1/0,-1,1/0,-1,-1} {
    \fill[carrier fill] (\x,\y,\z) circle (1.8pt);
  }
  \node[carrier site] at (0,0,0) {};
  \draw[carrier object] (1,1,0)--(1,0,1)--(0,1,1)--(-1,1,0)--(-1,0,1)--(0,-1,1)--(1,-1,0)--(1,0,-1)--(0,1,-1)--(-1,1,0);
  \draw[carrier muted, dashed] (1,1,0)--(0,1,1)--(-1,1,0);
  \fill[carrier hifill, opacity=0.5] (1,1,0)--(1,0,1)--(0,1,1)--cycle;
\end{tikzpicture}

  \caption{Кубооктаэдр --- выпуклая оболочка двенадцати соседей FCC.}
  \label{fig:carrier-cuboctahedron}
\end{figure}

\begin{figure}[htbp]
  \centering
  \begin{tikzpicture}[carrier panel]
  \draw[carrier muted] (-0.05,0) -- (1.3,0);
  \node[carrier site] at (0,0) {};
  \draw[carrier object] (0.72,0) -- (0.72,-0.38) -- (0.72,0.38);
  \draw[carrier arrow] (0,0) -- (0.72,0);
  \node[anchor=south, font=\scriptsize] at (0.36,-0.12) {$R_{\mathrm{in}}^{\mathrm{hull}}$};
  \draw[carrier arrow] (0,0) -- (0.55,0.48);
  \node[anchor=west, font=\scriptsize] at (0.6,0.5) {$R_{\mathrm{out}}=a$};
  \draw[carrier muted, dashed] (0.48,0.62) -- (0.72,0.62);
  \node[carrier muted, font=\scriptsize] at (0.6,0.75) {$a\sqrt{2/3}$};
\end{tikzpicture}

  \caption{Меридиан: квадратная и треугольная грани на разных расстояниях от центра.}
  \label{fig:carrier-cubocta-meridian}
\end{figure}

Объём ячейки Вороного считается отдельно.
Кубическая ячейка со стороной $A=a\sqrt{2}$ имеет объём $A^3=2\sqrt{2}\,a^3$ и содержит четыре узла FCC, поэтому на один узел
\[
  v_{\hv}=A^3/4=a^3/\sqrt{2}.
\]
Стена Вороного проходит посредине ребра к соседу, на расстоянии $a/2$ от узла.
Это не радиус оболочки:
\[
  R_{\mathrm{in}}^{\mathrm{Voronoi}}=a/2=\kgeom\, R_{\mathrm{in}}^{\mathrm{hull}}.
\]
(\cref{fig:carrier-hull-voronoi,fig:carrier-voronoi-cell}).
Двенадцать соседей двойственны двенадцати ромбам этой ячейки.

\begin{figure}[htbp]
  \centering
  \begin{tikzpicture}[carrier panel, scale=1.45]
  \node[carrier site] at (0,0) {};
  \draw[carrier object] (0,0) circle (1);
  \draw[carrier muted] (0,0) circle ({1/sqrt(2)});
  \draw[carrier muted, dashed] (0,0) circle (0.5);
  \draw[carrier arrow] (0,0) -- (1,0);
  \node[anchor=south, inner sep=1pt] at (0.5,-0.06) {$R_{\mathrm{out}}=a$};
  \draw[carrier arrow, carrier muted] (0,0) -- ({0.5/sqrt(2)},{0.5/sqrt(2)});
  \node[carrier muted, font=\scriptsize, anchor=west] at (0.42,0.38) {$R_{\mathrm{in}}^{\mathrm{hull}}$};
  \draw[carrier arrow, carrier muted, dashed] (0,0) -- (0.5,0);
  \node[carrier muted, font=\scriptsize, anchor=north] at (0.5,0.08) {$R_{\mathrm{in}}^{\mathrm{Voronoi}}$};
  \node[anchor=north, font=\footnotesize] at (0,-1.25)
    {$\kappa_{\mathrm{FCC}}=1/\sqrt{2}$};
\end{tikzpicture}

  \caption{Три радиуса однотактового тела на одном узле.}
  \label{fig:carrier-hull-voronoi}
\end{figure}

\begin{figure}[htbp]
  \centering
  \begin{tikzpicture}[carrier panel, iso view, scale=1.0]
  % schematic rhombic dodecahedron neighborhood
  \node[carrier site] at (0,0,0) {};
  \fill[carrier hifill, opacity=0.55]
    (1,0,0) -- (0,1,0) -- (-1,0,0) -- (0,-1,0) -- cycle;
  \draw[carrier object]
    (1,0,0) -- (0,1,0) -- (-1,0,0) -- (0,-1,0) -- cycle;
  \foreach \x/\y/\z in {1,1,0/1,-1,0/-1,1,0/-1,-1,0/1,0,1/1,0,-1/-1,0,1/-1,0,-1/0,1,1/0,1,-1/0,-1,1/0,-1,-1} {
    \draw[carrier muted, line width=0.4pt] (0,0,0) -- (\x*0.55,\y*0.55,\z*0.55);
    \fill[carrier fill, opacity=0.35] (\x*0.55,\y*0.55,\z*0.55) circle (1.5pt);
  }
  \draw[carrier object] (0.55,0.55,0) -- (0.55,0,0.55) -- (0,0.55,0.55) -- cycle;
\end{tikzpicture}

  \caption{Замощение $\mathbb{R}^3$ ячейками Вороного (ромбододекаэдры).}
  \label{fig:carrier-voronoi-cell}
\end{figure}

Объём самой оболочки находится вычитанием.
Многогранник с вершинами всех перестановок $(\pm 1,\pm 1,0)$ есть пересечение октаэдра $|x|+|y|+|z|\le 2$ и куба $|x|,|y|,|z|\le 1$.
Объём октаэдра $|x|+|y|+|z|\le 1$ равен $4/3$, поэтому при правой части~$2$ объём равен $(4/3)\cdot 8=32/3$.
Шесть шапок за гранями куба, по одной на полуось, имеют объём $2/3$ каждая: при $x$ от~$1$ до~$2$ сечение $|y|+|z|\le 2-x$ есть квадрат в повёрнутых осях площади $2(2-x)^2$, и
\[
  \int_1^2 2(2-x)^2\,\mathrm{d}x=\frac{2}{3}.
\]
Шесть шапок дают~$4$, остаётся $32/3-4=20/3$.
Ребро этого многогранника равно $\sqrt{2}$.
Масштаб к ребру~$a$ есть множитель $a/\sqrt{2}$, объём умножается на $(a/\sqrt{2})^3$:
\[
  V_{\mathrm{hull}}=\frac{20}{3}\cdot\frac{a^3}{2\sqrt{2}}=\frac{5\sqrt{2}}{3}\,a^3,
  \qquad
  \frac{V_{\mathrm{hull}}}{v_{\hv}}=\frac{5\sqrt{2}}{3}\cdot\sqrt{2}=\frac{10}{3}.
\]
Формула $(8\sqrt{2}/3)\,a^3$ с отношением $16/3$ --- объём другого тела.

Макроскопическая скорость связана с тактовой тем же отношением радиусов оболочки:
\[
  c=\kgeom c_0,\qquad \htick=\kgeom\, t_P,\qquad c_0=\sqrt{2}\,c,
\]
где $t_P=\hl/c$ --- учебниковая метка.
Непрерывное определение $t_P=\hl/c$ предполагало бы $c_0=c$; на этой оболочке тактовая и макроскопическая скорости различаются множителем~$\sqrt{2}$.

\section{Планковская решётка}
\label{sec:planck-lattice}

\begin{definition}[Планковская решётка]
\label{def:planck-lattice}
Носитель $\Lambda$ --- FCC-подрешётка $\mathbb{R}^3$
(теорема~\ref{thm:fcc-choice}).
Срез плоскостью $\{111\}$ даёт гексагональный слой $(2{+}1)$ (\cref{sec:hex-slice}, \cref{fig:carrier-fcc-111-slice}).
Объём элементарной ячейки Вороного
\[
  v_{\hv}=\hl^3/\sqrt{2}.
\]
\end{definition}

\begin{figure}[htbp]
  \centering
  \begin{tikzpicture}[carrier panel, iso view, scale=1.0]
  \def\r{0.28}
  \IsoSphereInk{0}{0}{0}{\r}
  \foreach \x/\y/\z in {
    1,1,0/1,-1,0/-1,1,0/-1,-1,0,
    1,0,1/1,0,-1/-1,0,1/-1,0,-1,
    0,1,1/0,1,-1/0,-1,1/0,-1,-1
  } {
    \IsoSphereMuted{\x}{\y}{\z}{\r}
  }
  \draw[carrier muted, dashed, fill=gray!8, opacity=0.45]
    (-1.2,-1.2,1.2) -- (1.2,-1.2,0) -- (1.2,1.2,-0.4) -- (-1.2,1.2,0.4) -- cycle;
  \node[carrier muted, font=\scriptsize, anchor=south west] at (-1.0,0.5,0.5) {$\{111\}$};
\end{tikzpicture}

  \caption{FCC-носитель и срез плоскостью $\{111\}$: вложенный гексагональный слой.}
  \label{fig:carrier-fcc-111-slice}
\end{figure}

Для этого носителя $\htick=\kgeom\, t_P$ при $\kgeom=1/\sqrt{2}$,
где $t_P=\hl/c$ --- планковское время в учебниковой нормировке; тогда $c_0=\sqrt{2}\,c$.
Это не подгонка задним числом: непрерывное определение $t_P=\hl/c$ предполагает $c_0=c$,
тогда как на решётке с $\kgeom=1/\sqrt{2}$ тактовая и макроскопическая скорости
принципиально различаются.

\section{Две скорости и мост $c=\kgeom c_0$}
\label{sec:two-c}

На гексагональном срезе $(2{+}1)$ коэффициент
\[
\kgeom_{\mathrm{hex}}=\sqrt{3}/2
\]
(\cref{sec:hex-slice}); на FCC-носителе
\[
\kgeom_{\mathrm{FCC}}=1/\sqrt{2}
\]
(\cref{sec:fcc-kappa}).
Величина $\kgeom$ --- отношение радиусов вписанной и описанной сфер однотактового тела.
Тождество $c=\kgeom c_0$ служит проверкой согласованности макроописания
после фиксации $\htick=\kgeom t_P$ (\cref{fig:carrier-two-speeds-micro,fig:carrier-two-speeds-macro}).

\begin{figure}[htbp]
  \centering
  \begin{tikzpicture}[carrier panel, x=2.1cm, y=2.1cm]
  \def\tlim{1.25}
  \def\xlim{1.35}
  \draw[carrier muted, dashed] (0,0) -- (\xlim,\tlim);
  \draw[carrier muted, dashed] (0,0) -- (-\xlim,\tlim);
  \draw[carrier object] (0,0) -- (\xlim,0) node[midway, below, inner sep=1pt] {$\hl$};
  \draw[carrier object] (0,0) -- (0,\tlim) node[midway, left, inner sep=1pt] {$\htick$};
  \foreach \i in {0,1,2,3} {
    \filldraw[carrier object] ({0.33*\i}, {0.33*\i}) circle (1.6pt);
  }
\end{tikzpicture}

  \caption{Микро: дискретный шаг по решётке, тактовая скорость $c_0$.}
  \label{fig:carrier-two-speeds-micro}
\end{figure}

\begin{figure}[htbp]
  \centering
  \begin{tikzpicture}[carrier panel, x=2.1cm, y=2.1cm]
  \def\tlim{1.25}
  \def\xlim{1.35}
  \def\czero{1.08}
  \def\kappa{0.7071}
  \pgfmathsetmacro{\cmacro}{\czero*\kappa}
  \draw[carrier muted, dashed] (0,0) -- (\xlim,\tlim);
  \draw[carrier muted, dashed] (0,0) -- (-\xlim,\tlim);
  \draw[carrier muted, dashed, line width=0.7pt] (0,0) -- (\xlim,\xlim/\czero);
  \draw[carrier object] (0,0) -- (\xlim,\xlim/\cmacro);
  \node[anchor=west, carrier muted, inner sep=1pt] at (0.55,0.62) {$c_0$};
  \node[anchor=west, inner sep=1pt] at (0.55,0.88) {$c=\kappa c_0$};
  \node[anchor=north west, carrier muted, font=\scriptsize, inner sep=1pt] at (0.05,0.95)
    {$\kappa=1/\sqrt{2}$};
\end{tikzpicture}

  \caption{Макро: осреднённый фронт, $c=\kgeom c_0$.}
  \label{fig:carrier-two-speeds-macro}
\end{figure}

\section{Отказ от континуума на микроуровне}
\label{sec:no-continuum-carrier}

Непрерывное пространство-время $(x,t)\in\mathbb{R}^4$ с пределом $\hl,\htick\to 0$
\emph{не} описывает микроуровень~$M$, как только зафиксированы FCC, $\kgeom=1/\sqrt{2}$
и значение~$c$ из эксперимента.
Непрерывные поля и дифференциальные уравнения --- язык макроуровня~$T$
на \emph{фиксированной} решётке; уменьшение $\hl$ не определено:
$\hl$ --- минимальный квант длины.

\section{Каузальность и наблюдаемая скорость света}

По решётке сигнал проходит за такт не дальше~$c_0\htick=\hl$.
На микроуровне мировая линия --- ломаная по рёбрам $N(x)$ со скоростью звена~$c_0$;
массивная частица --- зигзаг внутри того же конуса (\cref{fig:carrier-worldlines-fcc,fig:carrier-worldlines-hex}).
Макроскопический наблюдатель, осредняя поле на масштабах $\gg\hl$,
восстанавливает null-линию со скоростью $c=\kgeom c_0$ на FCC.
Локальная структура размера порядка~$\hl$ несёт конус~$c_0$;
составная структура, простирающаяся на многие ячейки, усредняется до макроскопической~$c$.
Сверхсветовое распространение на~$T$ здесь не возникает: сигнал не покидает
микроскопический световой конус~$c_0$.

\begin{figure}[htbp]
  \centering
  \begin{tikzpicture}[carrier panel, x=2.05cm, y=2.05cm]
  \def\tlim{1.2}
  \def\xlim{1.25}
  \def\kappa{0.7071}
  \draw[carrier muted, dashed] (0,0) -- (\xlim,\tlim);
  \draw[carrier muted, dashed] (0,0) -- (-\xlim,\tlim);
  \draw[carrier muted, line width=0.65pt]
    (0,0) -- (0.22,0.18) -- (0.38,0.08) -- (0.55,0.22) -- (0.72,0.12);
  \pgfmathsetmacro{\cmacro}{1.08/\kappa}
  \draw[carrier object] (0,0) -- (\xlim,\xlim/\cmacro);
  \node[anchor=west, font=\scriptsize, inner sep=1pt] at (0.45,0.35) {$c_0$};
  \node[anchor=west, font=\scriptsize, inner sep=1pt] at (0.45,0.72) {$c=\kappa c_0$};
\end{tikzpicture}

  \caption{Мировые линии на FCC ($\kgeom=1/\sqrt{2}$): микро-null, зигзаг и макро-фотон $c=\kgeom c_0$.}
  \label{fig:carrier-worldlines-fcc}
\end{figure}

\begin{figure}[htbp]
  \centering
  \input{figures/tikz/carrier-worldlines-hex.tex}
  \caption{Мировые линии на гекс-срезе $(2{+}1)$ ($\kgeom_{\mathrm{hex}}=\sqrt{3}/2$).}
  \label{fig:carrier-worldlines-hex}
\end{figure}

\section{Аналогия с физикой твёрдого тела}

В классической физике твёрдого тела атомы расположены на решётке Браве,
элементарная ячейка задаётся зоной Вигнера--Зейтца, колебания описываются фононами,
спектр --- зонами Блоха.
Если под атомами лежит среда элементарных носителей масштаба~$\hl$,
тот же словарь повторяется на более глубоком уровне:
центры --- ячейки~$\hv$ масштаба~$\ell_P$,
частицы --- первичные вихри и моды поля~$z$,
а вещество --- дефекты и составные узлы на вакуумной решётке.
Фононы вещества образуют второй кристалл над кипящим первым.

\section{Связь $(3{+}1)$ и $(2{+}1)$}

Гексагональный слой $(2{+}1)$ --- срез пространства Минковского $(3{+}1)$ с FCC
(\cref{fig:carrier-slice-hex-layer,fig:carrier-slice-fcc-layer}).
Мост $c=\kgeom c_0$ с $\kgeom=1/\sqrt{2}$ якорится на полном трёхмерном носителе;
чистый $(2{+}1)$ без объёма несёт $\kgeom=\sqrt{3}/2$.

\begin{figure}[htbp]
  \centering
  \begin{tikzpicture}[carrier panel]
  \CarrierHexRing{0.38}{0.65}
  \filldraw[carrier object] (0,0) circle (2pt);
\end{tikzpicture}

  \caption{Гексагональный срез $(2{+}1)$: $\kgeom_{\mathrm{hex}}=\sqrt{3}/2$, $|N|=6$.}
  \label{fig:carrier-slice-hex-layer}
\end{figure}

\begin{figure}[htbp]
  \centering
  \begin{tikzpicture}[carrier panel]
  \draw[carrier object] (0,0) circle (0.95);
  \foreach \a in {0,72,...,288} {
    \filldraw[carrier object, carrier fill] ({0.95*cos(\a)},{0.95*sin(\a)}) circle (2pt);
  }
  \filldraw[carrier object] (0,0) circle (2.5pt);
  \draw[carrier muted, dashed, line width=0.7pt] (-1.15,-0.55) -- (1.15,0.55);
  \node[anchor=south west, font=\scriptsize, carrier muted] at (-1.0,0.45) {$\{111\}$};
\end{tikzpicture}

  \caption{Носитель $(3{+}1)$ FCC и срез $\{111\}$: $\kgeom=1/\sqrt{2}$, $|N|=12$.}
  \label{fig:carrier-slice-fcc-layer}
\end{figure}

\section{Соседство среза и объёма}

На гексагональном срезе $|N(x)|=6$, на FCC $|N(x)|=12$.
Срез несёт шесть соседей и связь с объёмом.

